\Askhsh{33998}{4}
\begin{ylh}{1}
    \item 1.7 Όρια συνάρτησης στο άπειρο
    \item 3.3 Ορισμένο Ολοκλήρωμα
    \item 3.5 Η συνάρτηση $\mathrm{F}(x) = \int_\alpha^x f(t) dt$
\end{ylh}
\Askhseis{ΘΕΜΑ 4}

Το καπάκι ενός πεντάλιτρου δοχείου βενζίνης αφήνεται ανοιχτό τη χρονική στιγμή $t = 0$. Η βενζίνη που απομένει μέσα στο δοχείο συναρτήσει του χρόνου $t$ (σε εβδομάδες) δίνεται από τη συνεχή συνάρτηση $g(t)$ (σε λίτρα).

\begin{alist}
\item Να υπολογίσετε το ολοκλήρωμα $\int_0^2 5 \cdot \left(\frac{4}{5}\right)^t \cdot \ln \frac{4}{5} dt$. \monades{06}
\item Αν η βενζίνη του δοχείου έχει ρυθμό εξάτμισης που δίνεται από τον τύπο $g'(t)=5 \cdot \left(\frac{4}{5}\right)^t \cdot \ln \frac{4}{5}$, για κάθε $t > 0$, τότε να βρείτε τον όγκο της βενζίνης που περιέχει το δοχείο δύο εβδομάδες μετά το άνοιγμα του καπακιού του δοχείου. \monades{12}
\item Αν επιπλέον είναι γνωστό ότι η συνάρτηση που δίνει την ποσότητα της βενζίνης στο δοχείο μετά από $t$ εβδομάδες είναι η $g(t)=5 \cdot \left(\frac{4}{5}\right)^t$, $t \in [0, +\infty)$ τότε να διαπιστώσετε ότι καθώς ο χρόνος αυξάνεται απεριόριστα μόνο η μυρωδιά της βενζίνης θα υπάρχει στο δοχείο. \monades{07}
\end{alist}